\subsubsection{Componenti}
\label{sec:Componenti}

\paragraph{PageNavigation}
\textbf{Elementi chiave}:
\begin{itemize}
    \item \textbf{PreviousButton}: pulsante per navigare alla pagina precedente;
    \item \textbf{NextButton}: pulsante per navigare alla pagina successiva;
    \item \textbf{PageInput}: campo per inserire direttamente il numero di pagina;
    \item \textbf{PageIndicator}: mostra la pagina corrente e il numero totale;
\end{itemize}

\paragraph{SearchBar}
\textbf{Elementi chiave}:
\begin{itemize}
    \item \textbf{InputField}: campo di input per inserire parole chiave da cercare;
    \item \textbf{SearchIcon}: icona cliccabile per avviare la ricerca;
\end{itemize}

\paragraph{TestNameForm}
\textbf{Elementi chiave}:
\begin{itemize}
    \item \textbf{NewName}: campo di input per modificare il nome del test;
    \item \textbf{ConfirmButton}: pulsante per confermare il nuovo nome;
    \item \textbf{AbortButton}: pulsante per annullare la modifica;
    \item \textbf{ValidationFeedback}: messaggi di errore per input vuoti o duplicati;
\end{itemize}

\paragraph{TestListPage}
\textbf{Elementi chiave}:
\begin{itemize}
    \item \textbf{SearchBar}: barra per ricerca dei test salvati;
    \item \textbf{TestListContainer}: contenitore per la lista dei test salvati;
    \item \textbf{EmptyState}: messaggio e grafica quando non ci sono test disponibili;
\end{itemize}
\textbf{Servizio} \\
\textbf{TestService}: servizio iniettato tramite `@Injectable`, fornisce le operazioni per la gestione dell'elenco test.

\paragraph{TestListView}
\textbf{Elementi chiave}:
\begin{itemize}
    \item \textbf{TestItem}: elemento della lista dei test salvati;
    \item \textbf{UploadButton}: pulsante per caricare un test effettuato precedentemente; 
    \item \textbf{RenameButton}: pulsante per rinominare un test;
    \item \textbf{RemoveTestButton}: pulsante per rimuovere un test selezionato;
    \item \textbf{DoubleClickListener}: viene selezionato un test quando viene effettuato un doppio click su un elemento della lista;
\end{itemize}

\textbf{Input}:
\begin{itemize}
    \item \textbf{testList}: lista di test da visualizzare, fornito dal componente genitore;
\end{itemize}

\textbf{Output}:
\begin{itemize}
    \item \textbf{loadTest}: evento emesso al doppio click per richiedere il caricamento di un test;
    \item \textbf{deleteTest}: evento emesso quando si clicca il pulsante di rimozione;
    \item \textbf{renameTest}: evento per rinominare un test tramite `TestNameForm`;
\end{itemize}

\paragraph{TestPage}
\textbf{Elementi chiave}:
\begin{itemize}
    \item \textbf{TestIndex}: riepilogo sintetico delle metriche del test;
    \item \textbf{TestPageView}: sezione con risultati dettagliati del test;
    \item \textbf{ComparisonButton}: pulsante per avviare il confronto tra test;
    \item \textbf{SaveButton}: pulsante per salvare il test corrente;
\end{itemize}
\textbf{Servizio} \\
\textbf{TestService}: servizio iniettato tramite `@Injectable`, fornisce le funzionalità principali per la gestione dei test.

\paragraph{TestPageView}
\textbf{Elementi chiave}:
\begin{itemize}
    \item \textbf{ScatterDiagram}: visualizzazione grafica dei risultati tramite diagramma a dispersione;
    \item \textbf{CakeDiagram}: diagramma a torta che mostra la percentuale di risposte corrette;
    \item \textbf{BarDiagram}: diagramma a barre che mostra la distribuzione dei risultati per intervalli;
    \item \textbf{ResultsView}: contenitore dei risultati;
\end{itemize}

\textbf{Input}:
\begin{itemize}
    \item \textbf{Results}: array dei risultati da visualizzare;
\end{itemize}

\textbf{Output}:
\begin{itemize}
    \item \textbf{ChangePage}: evento emesso quando l'utente interagisce con `PageNavigation`;
\end{itemize}

\paragraph{CakeDiagram}
\textbf{Elementi chiave}:
\begin{itemize}
    \item \textbf{PieChart}: grafico a torta;
    \item \textbf{PieElement}: sezioni del grafico a torta;
    \item \textbf{Legend}: legenda che descrive i colori utilizzati;
    \item \textbf{Tooltip}: messaggio informativo al passaggio del mouse;
\end{itemize}

\textbf{Input}:
\begin{itemize}
    \item \textbf{CorrectnessValue}: percentuale di risposte corrette;
\end{itemize}

\paragraph{BarDiagram}
\textbf{Elementi chiave}:
\begin{itemize}
    \item \textbf{BarChart}: grafico a barre;
    \item \textbf{BarElement}: gruppi di barre per ogni categoria valutata (0-0.2, 0.2-0.4, 0.4-0.6, 0.6-0.8, 0.8-1);
    \item \textbf{BarLabel}: etichette per ogni barra;
    \item \textbf{BarTooltip}: messaggio informativo al passaggio del mouse;
\end{itemize}

\textbf{Input}:
\begin{itemize}
    \item \textbf{distribution}: array con il numero di risultati per intervallo;
\end{itemize}

\paragraph{ScatterDiagram}
\textbf{Elementi chiave}:
\begin{itemize}
    \item \textbf{Element}: distribuzione di punti XY per ogni risultato testato;
    \item \textbf{ZoomTool}: strumento per ingrandire una zona specifica del grafico;
\end{itemize}

\textbf{Input}:
\begin{itemize}
    \item \textbf{ElementValues}: coordinate dei punti da rappresentare;
\end{itemize}

\paragraph{TestComparisonPage}  
\textbf{Elementi chiave}:
\begin{itemize}
    \item \textbf{ComparisonIndex}: confronto sintetico tra due test;
    \item \textbf{ComparisonPageContainer}: contenitore della pagina di confronto;
    \item \textbf{BackButton}: pulsante per tornare alla test base usato per il confronto;
\end{itemize}
\textbf{Servizio} \\
\textbf{TestService}: servizio iniettato tramite `@Injectable`, espone i metodi per il confronto tra test.

\paragraph{ComparisonIndex}  
\textbf{Elementi chiave}:
\begin{itemize}
    \item \textbf{MetricComparison}: confronto tra le metriche chiave dei due test;
\end{itemize}

\textbf{Input}:
\begin{itemize}
    \item \textbf{TestIndex1}: metriche del primo test;
    \item \textbf{TestIndex2}: metriche del secondo test;
\end{itemize}

\paragraph{TestIndex}
\textbf{Elementi chiave}:
\begin{itemize}
    \item \textbf{DataSummary}: indice riassuntivo del test;
\end{itemize}

\textbf{Input}:
\begin{itemize}
    \item \textbf{summaryData}: oggetto contenente i valori sintetici del test;
\end{itemize}

\paragraph{ComparisonPageView}  
\textbf{Elementi chiave}:
\begin{itemize}
    \item \textbf{ComparisonScatter}: diagramma XY comparativo tra i due test;
    \item \textbf{DiffList}: lista dettagliata delle differenze tra risposte;
\end{itemize}

\textbf{Input}:
\begin{itemize}
    \item \textbf{comparisonData}: oggetto contenente le differenze per singola domanda;
\end{itemize}